\documentclass[11pt]{article}
\usepackage{amsmath, amssymb, amsthm, mathtools, bm}
\usepackage[margin=1.1in]{geometry}
\newtheorem{theorem}{Theorem}
\newtheorem{lemma}[theorem]{Lemma}
\newtheorem{proposition}[theorem]{Proposition}
\newtheorem{corollary}[theorem]{Corollary}
\newtheorem{assumption}{Assumption}
\newtheorem{definition}{Definition}
\DeclareMathOperator{\E}{\mathbb{E}}
\DeclareMathOperator{\Var}{Var}
\DeclareMathOperator{\argmin}{argmin}
\DeclareMathOperator{\atom}{atom}
\newcommand{\ind}{\mathbb{1}}
\newcommand{\indep}{\perp\!\!\!\perp}

\title{RuleOPE: Theoretical guarantees}
\author{}
\date{}

\begin{document}
\maketitle

\section{Setup and notation}

We observe i.i.d.\ tuples $(X_i, A_i, R_i, C_i)_{i=1}^N$ drawn from a
distribution $P$ on $\mathcal{X} \times \mathcal{A} \times [0, 1] \times
\{0, 1\}$, where $X_i \in \mathcal{X}$ encodes query and retrieval
features, $A_i \in \mathcal{A} = \{a_0, a_1, \ldots, a_K\}$ is the
logged action under a known logging policy $\pi_0$, $R_i \in [0, 1]$
is the logged reward, and $C_i \in \{0, 1\}$ is a sparse correction
signal.  The logging policy is $\pi_0(a \mid x) > 0$ for all $a, x$
(see Assumption~\ref{ass:positivity} below).

A \emph{rule} is a pair $\rho = (\phi, a_\rho)$ where
$\phi: \mathcal{X} \to \{0, 1\}$ is a Boolean predicate expressed as a
conjunction $\phi(x) = \prod_{\alpha \in S_\rho} \phi_\alpha(x)$ of
atomic predicates $\phi_\alpha$ from a finite vocabulary
$\mathcal{V}$, and $a_\rho \in \mathcal{A}$ is the action the rule
prescribes when $\phi(x) = 1$.  The \emph{policy induced by $\rho$} is
the deterministic function $\pi_\rho(x) = a_\rho$ if $\phi(x) = 1$ and
$\pi_\rho(x) = a_0$ otherwise, where $a_0$ denotes the default
(``noop'') action.

Let $R(x, a)$ denote the potential outcome (Rubin 1974) under
action $a$, so that $R_i = R(X_i, A_i)$.  The target of estimation is
\begin{equation}
V(\rho) \;=\; \E\left[R\bigl(X, \pi_\rho(X)\bigr)\right].
\label{eq:target}
\end{equation}

We will estimate $V(\rho)$ from the observed tuples for each rule
$\rho$ in a candidate set $\mathcal{R}$.

\section{Assumptions}

\begin{assumption}[Consistency of potential outcomes]\label{ass:consistency}
$R_i = R(X_i, A_i)$.
\end{assumption}

\begin{assumption}[Positivity]\label{ass:positivity}
There exists $\epsilon > 0$ such that $\pi_0(a \mid x) \geq \epsilon$
for all $a \in \mathcal{A}$ and $P$-a.s.\ $x$.
\end{assumption}

\begin{assumption}[Action unconfoundedness]\label{ass:action-uncon}
$\{R(x, a) : a \in \mathcal{A}\} \indep A \mid X$.
\end{assumption}

\begin{assumption}[Correction unconfoundedness]\label{ass:corr-uncon}
There exists a measurable $g: \mathcal{X} \times \mathcal{A} \to [0,1]$
such that $C \indep R \mid X, A$ and $P(C = 1 \mid X, A) = g(X, A)$.
\end{assumption}

Assumption~\ref{ass:consistency} is the standard SUTVA condition;
~\ref{ass:action-uncon} is the standard ignorability condition for
OPE;~\ref{ass:positivity} is required to inverse-weight observations.
Assumption~\ref{ass:corr-uncon} is new here and is the specific
assumption our correction-driven estimator relies on.  It says that
the correction indicator, conditional on features and logged action,
is independent of the potential reward---the expert's decision to
review does not depend on the specific reward realisation beyond what
is encoded in $(X, A)$.  We characterise below three RAG-specific
failure modes of this assumption.

\section{Compositional reward regression}

Let $\Phi(x) = (\phi_\alpha(x))_{\alpha \in \mathcal{V}} \in \{0, 1\}^d$
be the atomic indicator vector.  We model the conditional mean reward
as
\begin{equation}
m(x, a) \;=\; \E[R(X, a) \mid X = x] \;=\; \beta_{0, a} + \Phi(x)^\top \beta_a,
\label{eq:reg}
\end{equation}
with $\beta_a \in \mathbb{R}^d$.  The key observation: for a rule
$\rho$ with atom set $S_\rho$,
\begin{equation*}
\E[m(X, a_\rho) \mid \phi(X) = 1] = \beta_{0, a_\rho}
  + \sum_{\alpha \in S_\rho} \beta_{\alpha, a_\rho}
  + \sum_{\alpha \notin S_\rho} \beta_{\alpha, a_\rho}\, \E[\phi_\alpha(X) \mid \phi(X) = 1].
\end{equation*}
Two rules $\rho, \rho'$ that share atom $\alpha \in S_\rho \cap S_{\rho'}$
\emph{share the coefficient} $\beta_{\alpha, a_\rho}$ when $a_\rho =
a_{\rho'}$.  This is the source of the variance reduction proved in
Theorem~\ref{thm:variance}.

\section{The RuleOPE estimator}

Given cross-fitted estimators $\widehat m, \widehat g$ of
$m, g$ respectively, and using the shorthand
$w_i(\rho) = \ind\{A_i = \pi_\rho(X_i)\}/\pi_0(A_i \mid X_i)$, define
\begin{align}
\psi_i(\rho) \;&=\; \widehat m\bigl(X_i, \pi_\rho(X_i)\bigr) \nonumber \\
  &\quad + w_i(\rho)\bigl(R_i - \widehat m(X_i, A_i)\bigr) \nonumber \\
  &\quad + C_i\, h_i(\rho)\,\bigl(\tilde r_i(\rho) - \widehat m(X_i, \pi_\rho(X_i))\bigr),
\label{eq:rope}
\end{align}
and $\widehat V^{\text{ROPE}}(\rho) = N^{-1}\sum_i \psi_i(\rho)$.

The gate $h_i(\rho) = \ind\{\pi_\rho(X_i)\ne A_i\}\cdot h^\star(X_i, \rho)$
is a learnt correction-informativeness weight, and $\tilde r_i(\rho)$
is an imputation of the counterfactual reward.  For the action
$a_\rho = \texttt{abstain}$ we set $\tilde r_i \equiv r_{\rm abs}$;
for $a_\rho \in \{\texttt{filter}, \texttt{rerank}\}$ we set
$\tilde r_i = \widehat m(X_i, a_\rho)$.

\section{Consistency}

\begin{theorem}[Consistency]\label{thm:consistency}
Assume \ref{ass:consistency}--\ref{ass:corr-uncon}, and that either
\begin{itemize}
\item[(i)] $\widehat m \xrightarrow{P} m$, or
\item[(ii)] the effective propensity
  $\bar\pi_0(a_\rho \mid x) = \pi_0(a_\rho \mid x) +
  g(x, a_0)\,h^\star(x, \rho)\,\ind\{\phi(x) = 1\}$ is known exactly
  and $\widehat g \xrightarrow{P} g$.
\end{itemize}
Then $\widehat V^{\text{ROPE}}(\rho) \xrightarrow{P} V(\rho)$.
\end{theorem}

\begin{proof}[Proof sketch]
Under~\ref{ass:action-uncon},
$\E[w_i(\rho)(R_i - m(X_i, A_i))] = 0$, so the second term in
\eqref{eq:rope} has zero population mean and is a variance-reducing
correction centred at the truth.  Under~\ref{ass:corr-uncon},
$\E[C_i \mid X_i, A_i] = g(X_i, A_i)$ is independent of $R_i$, so
the third term likewise has zero population mean \emph{when} the gate
$h^\star$ is chosen so that
\[
\E\Bigl[C_i\,h^\star(X_i, \rho)\bigl(\tilde r_i(\rho) - m(X_i, \pi_\rho(X_i))\bigr)\Bigr]
= \E\bigl[\bigl(R(X, \pi_\rho(X)) - m(X, \pi_\rho(X))\bigr) \ind\{\phi(X)=1\}\bigr].
\]
This equality is a moment condition that can be solved for $h^\star$ in
closed form under linear regression (eq.~\ref{eq:reg}); see
Appendix~A.  Standard DR arguments (Robins et al.\ 1994,
Chernozhukov et al.\ 2018) then give either (i) or (ii) as a
sufficient condition for consistency.
\end{proof}

\section{Variance reduction from compositional regression}

\begin{theorem}[Compositional variance bound]\label{thm:variance}
Suppose the reward model~\eqref{eq:reg} is well-specified, that
$\|\Phi(X)\|_2 \le B$ a.s., and that the ridge regression is trained
with regularisation $\lambda > 0$.  Let $\mathcal{R}$ be a collection
of $M$ rules each of depth at most $D$, with atom-vocabulary
$\mathcal{V}$ of size $d$, and let $\mathcal R$ use actions in
$\{a_0, \ldots, a_K\}$.  Then the \emph{total} estimator variance
across $\mathcal{R}$ satisfies
\begin{equation*}
\sum_{\rho \in \mathcal{R}} \Var\!\bigl(\widehat V^{\text{ROPE}}(\rho)\bigr)
\;\le\;
\frac{1}{N}\Bigl(\sigma_R^2 \cdot \max_{i,\rho} w_i(\rho)^2 \cdot M
+ \underbrace{K (d + 1) \cdot \kappa(\lambda)}_{\textrm{shared regression}}\Bigr),
\end{equation*}
where $\kappa(\lambda) = O(B^2/\lambda)$ is the regression's effective
complexity and $\sigma_R^2$ is the conditional reward variance.
\end{theorem}

\begin{proof}[Proof sketch]
The first term is the standard IPS variance $\times$ rule count, which
is unavoidable.  The second term follows because the regression
parameters $(\beta_{0, a}, \beta_{\cdot, a})$ are shared across all
rules using action $a$, so the regression's contribution to
$\sum_\rho \Var$ is bounded by the regression's own parameter
complexity $K(d + 1)$ times an effective-degrees-of-freedom
$\kappa(\lambda)$, independent of $M$.  Thus when $M \gg d$, the
regression term is \emph{sublinear} in $M$ compared to the $O(M)$
cost of per-rule regression.
\end{proof}

The practical consequence: doubling the rule set doubles only the IPS
part of the variance, not the regression part.  Baselines that fit a
separate regression per rule scale linearly with $M$ in both parts.

\section{Failure modes of Assumption~\ref{ass:corr-uncon}}

\paragraph{Query-dependent correction effort.}
If $P(C=1 \mid X, A)$ depends on a feature that is \emph{also} a
direct parent of $R$ not captured by $\Phi(X)$, the gate term biases
the estimator.  Diagnostic: $\widehat g(x, a_0)$ has residual
correlation with $R \mid X, A$ on held-out data.

\paragraph{Self-consistent answer bias.}
When the generator's confidence $\operatorname{gen\_conf}(X)$ causes
experts to \emph{under}-review confidently wrong answers,
Assumption~\ref{ass:corr-uncon} fails in a systematic direction.  The
estimator over-weights high-reward queries in the correction term.
Mitigation: augment $\Phi$ with a confidence atom (we include
$\operatorname{gen\_conf}<\tau$ atoms in $\mathcal{V}$).

\paragraph{Corpus drift.}
If the logging distribution of $X$ changes between training and
evaluation (train-test mismatch on $\Phi(X)$), positivity degrades and
the IPS term's variance blows up.  Detected by ESS $\ll N$ on
held-out logs; mitigated by reweighting via a covariate shift
estimator.

\section{Remarks}

The estimator satisfies standard DR consistency at root-$N$ rates
whenever the regression converges faster than $N^{-1/4}$ and the
propensity or correction model at the same rate (``double machine
learning'' regime).  For linear ridge regression in fixed dimension
$d = |\mathcal{V}|$ this is automatic.

%---------------------------------------------------------------------------
\section{Cross-rule and between-estimator shrinkage}

We introduce two shrinkage estimators on top of RuleOPE that provably
reduce estimation error when evaluating a \emph{set} of rules.

\subsection{Between-estimator shrinkage: DualShrinkOPE}

Fix a rule $\rho$.  Let $\widehat V_{\rm DR}(\rho), \widehat V_{\rm DM}(\rho)$
denote the RuleOPE and Direct Method point estimates, with respective
biases $b_{\rm DR}, b_{\rm DM}$ and variances $\sigma_{\rm DR}^2,
\sigma_{\rm DM}^2$.  For any weight $w \in [0, 1]$ define
$$
\widetilde V_w(\rho) = w\,\widehat V_{\rm DM}(\rho) + (1 - w)\,\widehat V_{\rm DR}(\rho).
$$
\begin{proposition}[Bayes-optimal weight]\label{prop:dual-w}
The MSE-minimising weight is
$$
w^\star(\rho) = \frac{\sigma_{\rm DR}^2 + b_{\rm DR}^2}{\sigma_{\rm DR}^2 + b_{\rm DR}^2 + \sigma_{\rm DM}^2 + b_{\rm DM}^2} \in [0, 1],
$$
and the resulting estimator satisfies
$$
\mathrm{MSE}(\widetilde V_{w^\star}) = \frac{(\sigma_{\rm DR}^2 + b_{\rm DR}^2)(\sigma_{\rm DM}^2 + b_{\rm DM}^2)}{\sigma_{\rm DR}^2 + b_{\rm DR}^2 + \sigma_{\rm DM}^2 + b_{\rm DM}^2}
\le \min\bigl(\sigma_{\rm DR}^2 + b_{\rm DR}^2,\, \sigma_{\rm DM}^2 + b_{\rm DM}^2\bigr).
$$
\end{proposition}
\begin{proof}
Direct computation and AM-GM.  $\Box$
\end{proof}

When $b_{\rm DR}$, $b_{\rm DM}$ are unknown we replace them by the
finite-sample proxy $\widehat b = |\widehat V_{\rm DR} - \widehat V_{\rm DM}|$
and plug into $w^\star$.  A rule-OPE analogue of Wang et al.'s switch
estimator, but with a soft, per-rule data-driven weight.  Section 7 of
the paper documents a $10$--$23\%$ MSE reduction over DR in the
deterministic-logging regime.

\subsection{Cross-rule random-effects shrinkage: JointRuleOPE}

\begin{theorem}[Compositional random-effects dominance]\label{thm:shrink}
Let $\{\widehat V(\rho_m)\}_{m=1}^M$ be per-rule estimates with
respective variances $\sigma_m^2$.  Suppose the rule values decompose as
$V(\rho_m) = \mu_m + \eta_m$ where $\mu_m = F(\rho_m)^\top \beta$ is
linear in rule features $F$ and $\eta_m \sim \mathcal{N}(0, \tau^2)$
i.i.d.  Then the random-effects estimator
$$
\widetilde V(\rho_m) = w_m\, \widehat\mu_m + (1 - w_m)\, \widehat V(\rho_m),\quad
w_m = \frac{\sigma_m^2}{\sigma_m^2 + \tau^2},
$$
has joint MSE strictly smaller than the independent estimator
whenever $\tau^2 < \infty$, and asymptotically achieves the lower
envelope
$$
\sum_m \mathrm{MSE}(\widetilde V) \le \sum_m \frac{\sigma_m^2 \tau^2}{\sigma_m^2 + \tau^2}
< \sum_m \sigma_m^2 = \sum_m \mathrm{MSE}\bigl(\widehat V(\rho_m)\bigr).
$$
\end{theorem}
\begin{proof}
Classical random-effects / Efron--Morris argument applied to the
regression-induced prior $\mu_m = F(\rho_m)^\top \beta$.  Cross-fitting
of $\widehat\mu$ is required for the plug-in estimator to inherit the
dominance in finite samples.  $\Box$
\end{proof}

The practical form replaces $\tau^2$ with the DerSimonian--Laird
estimator $\widehat\tau^2 = \max(0, \overline{r^2 - \sigma^2})$ where
$r_m = \widehat V(\rho_m) - \widehat\mu_m$, and uses the same
per-rule $w_m$ above.

%---------------------------------------------------------------------------
\section{Pessimistic rule selection with compositional complexity}

\subsection{Setup}

Given estimates $\{(\widehat V_m, \widehat\sigma_m)\}_{m=1}^M$ we
select a rule by maximising a lower confidence bound
$$
\widehat V_{\rm LCB}(\rho_m) = \widehat V_m - c_M \widehat\sigma_m.
$$
The regret against the oracle rule $\rho^\dagger =
\argmax_m V(\rho_m)$ is $V(\rho^\dagger) - V(\widehat\rho)$.

\subsection{Standard union-bound LCB}

\begin{proposition}\label{prop:lcb-union}
Under sub-Gaussian influence functions with scale $\sigma_m$, setting
$c_M = \sqrt{2 \log(M/\delta)}$ yields, with probability $1 - \delta$,
$\max_m |\widehat V_m - V_m| \le c_M \widehat\sigma_m$, and hence the
regret bound
$$
V(\rho^\dagger) - V(\widehat\rho) \le 2 c_M \max_m \widehat\sigma_m = O\!\left(\sqrt{\log(M/\delta)/N}\right).
$$
\end{proposition}

\subsection{Atom-sparse LCB (ours)}

Let the rule-value function be $s$-sparse in the atom-compositional
basis: $V(\rho) = F(\rho)^\top \beta^\star$ with
$\|\beta^\star\|_0 \le s$ and $\|F(\rho)\|_\infty \le 1$.  The
hypothesis class of candidate value functions is
$\mathcal{F}_s = \{F \mapsto F^\top \beta : \|\beta\|_0 \le s\}$.

\begin{theorem}[Compositional pessimistic regret]\label{thm:comp-lcb}
With $c_M = \sqrt{2 ((s + 1) \log(d + 1) + \log(1/\delta))}$, the
atom-sparse LCB selector has regret, with probability $1-\delta$,
$$
V(\rho^\dagger) - V(\widehat\rho) \le 2 c_M \max_m \widehat\sigma_m
= O\!\left(\sqrt{(s \log d + \log(1/\delta))/N}\right).
$$
\end{theorem}
\begin{proof}[Sketch]
The Rademacher complexity of $\mathcal{F}_s$ on $M$ rules is
$O(\sqrt{s \log d / N})$ by a standard Dudley chaining argument over
$s$-sparse linear functions.  Converting to an LCB exponent gives
$c_M = O(\sqrt{s \log d + \log(1/\delta)})$.  The regret inequality is
the usual LCB argument.  $\Box$
\end{proof}

\subsection{When does compositional dominate union-bound?}

Compare $s \log d + \log(1/\delta)$ to $\log M + \log(1/\delta)$.
The compositional bound is smaller whenever
$$
s < \log M / \log d.
$$
For our benchmark ($d = 48$, $M = 500$), this gives
$s < \log(500) / \log(48) \approx 1.6$.  The compositional bound
strictly dominates whenever the best rule depends on $\le 1$ atom
(depth-1 rules), and matches otherwise.  For a larger rule space
($M = 5000$ at depth 3), the threshold rises to $s < 2.2$, so
depth-2 rules with at most one truly informative atom benefit from
the compositional bound.

%---------------------------------------------------------------------------
\section{Counterfactual rule risk minimisation}

\begin{theorem}[CRRM regret]\label{thm:crrm}
Let $\widehat\rho = \argmax_\rho \bigl(\widehat V_{\rm LCB}(\rho) - \lambda\, |S_\rho|\bigr)$ with the
compositional LCB of Theorem~\ref{thm:comp-lcb}.  Assume A1--A4, the
base estimator has influence function $\psi$ bounded by $B$ a.s., and
the regression is well-specified.  Then with probability $1 - \delta$
over the training draw,
$$
V(\rho^\dagger) - V(\widehat\rho) \le 2\, \underbrace{\mathcal{R}_N(\mathcal{F}_s)}_{\textrm{Rademacher}} + 2B\sqrt{\frac{2 \log(1/\delta)}{N}} + \lambda\, |S_{\rho^\dagger}|,
$$
where $\mathcal{R}_N(\mathcal{F}_s) = O\bigl(\sqrt{s \log d / N}\bigr)$.
\end{theorem}
\begin{proof}
Standard Rademacher generalisation bound applied to the influence-function
class $\{\psi_\rho : \rho \in \mathcal{R}\}$.  Because the influence function
of RuleOPE is linear in $F(\rho)$, the Rademacher complexity is controlled
by $\mathcal{F}_s$.  $\Box$
\end{proof}

The CRRM regret scales with the atom sparsity $s$ and the log of the
atom vocabulary $d$, not with the (exponential-in-depth) rule-space
size $M$.  When $s$ is small, CRRM is strictly tighter than ERM with
union-bound selection.

%---------------------------------------------------------------------------
\section{Identification and efficiency under deterministic logging}
\label{sec:identification}

This section is the theoretical spine of the paper.  We prove that under
deterministic logging $A \equiv a_0$ the target $V(\rho)$ is
\emph{not point-identified} from $(X, A, R)$ alone under standard
ignorability (A1--A4), but becomes point-identified under an additional
structural assumption A5 on the correction model.  We derive the
corresponding efficient influence function (EIF) and show that RuleOPE
attains the semiparametric efficiency bound while estimators that ignore
the correction signal do not.

\subsection{Deterministic logging and the partitioned target}

Under deterministic logging $A_i \equiv a_0$ P-a.s.  Let
$p(x) = \mathbb{1}[\phi_\rho(x) = 1]$ be the rule's firing indicator and
$q(x) = 1 - p(x)$.  By linearity,
\begin{equation}
V(\rho) \;=\; \underbrace{\E[q(X) R(X, a_0)]}_{V_0(\rho)} \;+\; \underbrace{\E[p(X) R(X, a_\rho)]}_{V_1(\rho)}.
\label{eq:partition}
\end{equation}

\begin{lemma}[$V_0$ is trivially identified]
\label{lem:V0}
Under A1, A3, and deterministic logging, $V_0(\rho) = \E[q(X) R]$ is
identified from $P$, with unbiased plug-in $\widehat V_0 = N^{-1}\sum_i q(X_i) R_i$.
\end{lemma}
\begin{proof}
$R_i = R(X_i, A_i) = R(X_i, a_0)$ P-a.s., so $q(X_i) R_i = q(X_i) R(X_i, a_0)$.
Law of large numbers. $\Box$
\end{proof}

The hard object is $V_1(\rho)$, the \emph{counterfactual contribution
from firing records}.

\subsection{Theorem A: non-identifiability and sharp bounds under A1--A4}

\begin{theorem}[Sharp partial-identification bounds]
\label{thm:partial-id}
Assume A1--A4, $R \in [0, 1]$, and deterministic logging.  Then $V(\rho)$
is \emph{not} point-identified from $P$; the sharp identified interval
is
\begin{equation}
V_L(\rho) \;\le\; V(\rho) \;\le\; V_U(\rho), \qquad
\begin{aligned}
V_L(\rho) &= \E[q(X) R], \\
V_U(\rho) &= \E[q(X) R] + \E[p(X)].
\end{aligned}
\label{eq:sharp-bounds}
\end{equation}
The width of the identified interval is $\E[p(X)]$, the rule's firing probability.
\end{theorem}
\begin{proof}
\textbf{Sharpness of $V_U$.}  Consider the data-generating distribution
$P^\dagger$ obtained from $P$ by setting $R(x, a_\rho) \equiv 1$ P-a.s.
This choice does not conflict with any identified marginal of $P$:
observed $(X, R, C)$ are unchanged because $R_i = R(X_i, a_0)$ under
deterministic logging, and A4 pins $C \perp R | X, A$ so the correction
marginals $g(x, a_0)$ and $g(x, a_\rho)$ are unaffected by changes in
$R(\cdot, a_\rho)$.  Under $P^\dagger$, $V(\rho) = \E[q R(X, a_0)] + \E[p \cdot 1] = V_L(\rho) + \E[p]$.

\textbf{Sharpness of $V_L$.}  Symmetric: set $R(x, a_\rho) \equiv 0$.

\textbf{Non-identifiability.}  Both $P^\dagger$ values induce the same
observed-data marginal $P_{(X, A, R, C)}$, and yield different $V(\rho)$
whenever $\E[p(X)] > 0$.  $\Box$
\end{proof}

\begin{remark}
Theorem~\ref{thm:partial-id} says the identification gap is proportional
to the rule's firing probability $\E[p]$.  A rule that fires on 30\% of
queries has an identification gap of $0.3$ in the reward scale --- a
dominant source of error compared to any estimator's sampling variance.
No estimator using only $(X, A, R)$ can close this gap.
\end{remark}

\subsection{Theorem B: classical DR converges to a biased point inside the interval}

Let $\widehat V^{\rm DR}(\rho)$ denote the classical doubly-robust
estimator (equivalently CIPS-DR and CascadeDR all collapse to this under
deterministic logging, since the importance weight is zero on every
target-action record).  Let $m_\star(x, a) = \lim_N \widehat m(x, a)$
be the probability limit of the ridge regression.

\begin{theorem}[DR probability limit and bias]
\label{thm:dr-bias}
Under A1--A4, deterministic logging, and standard regularity conditions on
the ridge regression,
\begin{equation}
\widehat V^{\rm DR}(\rho) \;\xrightarrow{P}\; \E[q(X) R] + \E[p(X)\, m_\star(X, a_\rho)].
\label{eq:dr-limit}
\end{equation}
The bias relative to the true $V(\rho)$ is
\begin{equation}
\mathrm{Bias}(\widehat V^{\rm DR}(\rho)) \;=\; \E[p(X) \,(m_\star(X, a_\rho) - m(X, a_\rho))],
\label{eq:dr-bias}
\end{equation}
which is in general non-zero and equals zero only when the regression
is correctly specified at the unidentified counterfactual $m(\cdot, a_\rho)$.
\end{theorem}
\begin{proof}
Under deterministic logging, the importance-weight indicator
$\mathbb{1}[A_i = \pi_\rho(X_i)] = q(X_i)$, so
\begin{align*}
\widehat V^{\rm DR}(\rho)
&= \tfrac{1}{N}\sum_i \bigl[\widehat m(X_i, \pi_\rho(X_i)) + q(X_i) (R_i - \widehat m(X_i, a_0))\bigr] \\
&\xrightarrow{P} \E[p(X) m_\star(X, a_\rho)] + \E[q(X) m_\star(X, a_0)] + \E[q(X)(R - m_\star(X, a_0))] \\
&= \E[p(X) m_\star(X, a_\rho)] + \E[q(X) R] = V_0(\rho) + \E[p(X) m_\star(X, a_\rho)].
\end{align*}
Comparing to $V(\rho) = V_0 + \E[p \cdot m(X, a_\rho)]$ yields the bias. $\Box$
\end{proof}

\begin{corollary}[DR cannot close the identification gap]
Theorem~\ref{thm:dr-bias} shows the classical DR estimator converges to
a specific point in $[V_L, V_U]$ determined by the regression's
\emph{out-of-support extrapolation} at $a = a_\rho$.  Since the ridge
regression is trained only on data at $A = a_0$, its prediction at
$a = a_\rho$ is an unidentified extrapolation, and its asymptotic bias is
bounded only by the identification gap $\E[p(X)]$, not by any data-driven
quantity.  The same conclusion applies verbatim to CIPS-DR (the clip
eliminates any remaining importance-weight signal) and CascadeDR (the
position-wise factorisation still lives inside the same regression-only
family).
\end{corollary}

\subsection{A5: a bridge-function assumption (revised)}
\label{sec:a5-revised}

We now introduce a structural assumption on the correction mechanism
linking the observed distribution at $A = a_0$ to the counterfactual
reward regression at $A = a_\rho$.  The form mirrors the
\emph{proxy-variable / bridge-function} approach to proximal causal
identification (Miao, Geng, and Tchetgen Tchetgen 2018; Kallus, Mao,
and Uehara 2021) with the correction $C$ playing the role of the
negative-control exposure.

\paragraph{Erratum on an earlier version.}
A previous version of this manuscript defined the bridge as an
$\mathcal{X}$-measurable function $b_\rho: \mathcal{X} \to \mathbb{R}$
satisfying $\E[b_\rho(X)(C - g(X, a_0)) \mid X, A = a_0] = m(X, a_\rho) - m(X, a_0)$.
That equation is vacuous: for any $\mathcal{X}$-measurable $b_\rho$,
the LHS factors as $b_\rho(X)\,\E[C - g(X, a_0) \mid X, A = a_0] = b_\rho(X)\cdot 0 = 0$
identically, forcing $m(X, a_\rho) \equiv m(X, a_0)$.  The revised
definition below corrects this by allowing the bridge to depend on $C$
as well, which is the standard Miao--Geng--Tchetgen-Tchetgen bridge
form.  All downstream theorems (C, D, E, F) are restated with the
corrected definition; the substantive content (variance-reduction
bound, efficiency gap formula) carries through under a scope
narrowing from ``deterministic logging'' to ``deterministic logging
with an auxiliary pilot'' or, more commonly in this paper's
experiments, ``stochastic logging with correction-informed variance
reduction''.

\begin{definition}[Action bridge function, corrected]\label{def:bridge}
An \emph{action bridge function} for a rule $\rho$ is a measurable
map $b_\rho: \{0, 1\} \times \mathcal{X} \to \mathbb{R}$ satisfying,
for P-a.e.\ $x$,
\begin{equation}
\E\bigl[b_\rho(C, X) \,\big|\, X = x,\, A = a_0\bigr]
\;=\; m(x, a_\rho).
\label{eq:bridge}
\end{equation}
\end{definition}

The LHS is a functional of the observed distribution
$P_{(X, A=a_0, R, C)}$; the RHS is the counterfactual outcome
regression we need to identify $V_1(\rho)$.  Existence of $b_\rho$
transforms unobserved counterfactuals into observable functionals.
Expanding the conditional expectation at a point,
$\E[b_\rho(C, X) | X = x, A = a_0] = b_\rho(1, x)\,g(x, a_0) + b_\rho(0, x)(1 - g(x, a_0))$,
so $b_\rho$ is characterised by a pair of $\mathcal{X}$-measurable
functions $(b_\rho(0, \cdot), b_\rho(1, \cdot))$ satisfying one
linear functional equation at each $x$.  Uniqueness of $b_\rho$ is
not required for identification of $V(\rho)$; following the
proximal-identification literature, any solution to \eqref{eq:bridge}
gives the same value of $\E[p(X)\, b_\rho(C, X)]$.

\begin{assumption}[A5, bridge existence, revised]\label{ass:A5}
An action bridge function $b_\rho$ as in
Definition~\ref{def:bridge} exists for every rule $\rho$ of interest.
\end{assumption}

\paragraph{When does A5 hold?}
A5 is a non-trivial restriction.  The bridge equation~\eqref{eq:bridge}
is solvable for \emph{some} $b_\rho$ if and only if
$m(\cdot, a_\rho)$ lies in the closed range of the operator
$h \mapsto \E[h(C, \cdot) | A = a_0]$ acting on $L^2(P_{C, X | a_0})$.
Concretely, A5 is plausible when one of the following holds:
\begin{enumerate}
\item[(i)] \textbf{Stochastic logging with positivity on $a_\rho$.}
  The logging policy $\pi_0$ has $\pi_0(a_\rho | x) \ge \epsilon > 0$.
  Then $m(\cdot, a_\rho)$ is directly identified by A3 (regression on
  $(X, A, R)$), and \emph{any} $b_\rho$ matching the identified
  $m(\cdot, a_\rho)$ solves~\eqref{eq:bridge}; the bridge is a
  device for \emph{variance reduction} via the correction signal,
  not for identification.  This is the regime of all three real-data
  benchmarks in §7C.
\item[(ii)] \textbf{Deterministic logging plus pilot data on $a_\rho$.}
  A small on-policy exploration dataset $\mathcal{D}_{\rm pilot}$
  with $A = a_\rho$ is available.  $m(\cdot, a_\rho)$ is identified
  from $\mathcal{D}_{\rm pilot}$; the bridge is again a
  variance-reduction device.
\item[(iii)] \textbf{Parametric proxy structure.}  A parametric
  structural model (e.g.\ a \emph{double}-proxy Miao-style model with
  both an NCE and an NCO) yields point identification of
  $m(\cdot, a_\rho)$.  The correction-linearity sufficient form
  below is a simple instance, but is not sufficient by itself
  without one of (i)--(iii).
\end{enumerate}
Under A1--A4 and \emph{strictly} deterministic logging with no auxiliary
signal, Theorem~\ref{thm:partial-id} (non-identification, sharp
bounds) applies and A5 generically \emph{fails}; this is consistent
with the partial-identification structure rather than in conflict
with it.

\paragraph{Sufficient structural form (revised): correction-linearity
under scenario (i) or (ii).}  Consider the correction-linearity
model
\begin{equation}
g(x, a) \;=\; \alpha(x) + \beta(a) \cdot (1 - m(x, a)),\qquad \beta(a_0) \ne 0.
\label{eq:corr-linear}
\end{equation}
Under scenario~(i) with stochastic logging at $a_\rho$, or under
scenario~(ii) with pilot data, $m(x, a_0)$, $m(x, a_\rho)$,
$g(x, a_0)$, and the slopes $\beta(a_0), \beta(a_\rho)$ are all
identifiable from the observed + auxiliary data.  A natural
parameterisation of the bridge under correction-linearity is
\begin{equation}
b_\rho(c, x) \;=\; m(x, a_\rho) + \frac{\beta(a_\rho) - \beta(a_0)}{\beta(a_0)^2}\,(c - g(x, a_0)),
\label{eq:bridge-form}
\end{equation}
which satisfies \eqref{eq:bridge}:
$\E[b_\rho(C, X) | X = x, A = a_0] = m(x, a_\rho) + \frac{\beta(a_\rho) - \beta(a_0)}{\beta(a_0)^2}\cdot 0 = m(x, a_\rho)$,
and whose second term is the \emph{variance-reduction projection} of
the correction residual $(C - g(X, a_0))$ onto the counterfactual
contrast.  The scalar $(\beta(a_\rho) - \beta(a_0))/\beta(a_0)^2$
is the same quantity that appears in the constant bridge of the
earlier manuscript version; its interpretation is corrected to
``variance-reduction gain per unit correction-residual'' rather
than ``identification bridge''.

\begin{remark}[What correction-linearity buys]
Under scenario~(i), A3 alone identifies $V(\rho)$; A5 +
correction-linearity buys a \emph{strict variance improvement}
whenever $\beta(a_\rho) \ne \beta(a_0)$ and $g(x, a_0) \in (0, 1)$.
This is the content of Theorem~\ref{thm:inefficient} below.  Under
scenario~(ii), A5 + correction-linearity identifies $V(\rho)$ from
the union of pilot and main-logging data with variance strictly
smaller than the pilot-only plug-in.
\end{remark}

\subsection{Theorem C: point identification of V(ρ) (revised)}

\begin{theorem}[Point identification via bridge, revised]\label{thm:point-id}
Under A1, A3, and A5 (corrected),
\begin{equation}
V(\rho) \;=\; \E\bigl[q(X)\, R\bigr] \;+\; \E\bigl[p(X)\,\E[b_\rho(C, X) \mid X, A = a_0]\bigr],
\label{eq:identified-V}
\end{equation}
and each term on the RHS is identified from $P_{(X, A=a_0, R, C)}$
together with whichever auxiliary signal makes A5 non-vacuous
(§\ref{sec:a5-revised}, scenarios (i)--(iii)).  Under deterministic
logging at $A = a_0$ and scenario~(ii) with pilot data, the bridge
is estimated from the combined main+pilot sample and~\eqref{eq:identified-V}
reduces to a pilot-augmented regression plug-in with a
correction-driven variance-reduction term.
\end{theorem}
\begin{proof}
By A3, $\E[R | X = x, A = a_\rho] = m(x, a_\rho)$, so
$V_1(\rho) = \E[p(X) m(X, a_\rho)]$.
By A5, $\E[b_\rho(C, X) | X = x, A = a_0] = m(x, a_\rho)$ P-a.e., hence
$\E[p(X) m(X, a_\rho)] = \E[p(X)\,\E[b_\rho(C, X) | X, A = a_0]]$.
Combining with Lemma~\ref{lem:V0} gives~\eqref{eq:identified-V}.
Identification of the RHS: $q(X) R$ is observable; $g(X, a_0)$ and
$b_\rho$'s parameters are identifiable under scenarios (i) or (ii)
by construction. $\Box$
\end{proof}

\begin{remark}[Comparison with the earlier vacuous form]
The earlier version used the bridge residual $(C - g(X, a_0))$
multiplied by an $\mathcal{X}$-measurable $b_\rho(X)$.  That term
is \emph{identically zero} in conditional expectation and hence
cannot transport identification from $A = a_\rho$ to $A = a_0$.
The corrected \eqref{eq:bridge} is the standard bridge equation
($\E[b_\rho(C, X) | X, A = a_0] = m(X, a_\rho)$) which does
transport identification, at the cost of requiring $b_\rho$ to be
$(C, X)$-measurable and of narrowing the scope (scenarios
(i)--(iii)).  This correction brings the assumption into line with
the proximal-identification literature.
\end{remark}

\subsection{Theorem D: efficient influence function (revised)}

We derive the EIF of the identified functional \eqref{eq:identified-V}
under A1, A3, A5, working in scenario~(i) (stochastic logging) where
the full regression $m(x, a)$ is identifiable from $(X, A, R)$.  The
EIF under scenario~(ii) (deterministic logging + pilot) is analogous
with the pilot sample entering the logged-DR term.

\begin{theorem}[EIF, semiparametric efficiency bound, revised]\label{thm:eif}
Under A1, A3, A5 and stochastic logging with
$\pi_0(a | x) \ge \epsilon > 0$ for all $a$,
the efficient influence function of $V(\rho)$ at $P$ is
\begin{equation}
\psi^\star(O) \;=\;
\underbrace{m(X, \pi_\rho(X)) - V(\rho)}_{\mathrm{DM~plug\mbox{-}in}}
\,+\,
\underbrace{\frac{\mathbb{1}[A = \pi_\rho(X)]}{\pi_0(A | X)}\,(R - m(X, A))}_{\mathrm{DR~at~logged~action}}
\,+\,
\underbrace{p(X)\bigl(b_\rho(C, X) - \E[b_\rho(C, X) | X, A = a_0]\bigr)}_{\mathrm{bridge~variance\mbox{-}reduction}}.
\label{eq:eif}
\end{equation}
The semiparametric efficiency bound is
\begin{equation}
\mathrm{SEB}(\rho) \;=\; \Var_P(\psi^\star(O)).
\label{eq:seb}
\end{equation}
\end{theorem}
\begin{proof}[Sketch]
\textbf{(a) EIF under A1, A3 without A5.}  The target functional
$V(\rho) = \E[m(X, \pi_\rho(X))]$ under stochastic logging has the
classical DR EIF given by the first two terms of \eqref{eq:eif}
(Robins, Rotnitzky, Zhao 1994; Bickel et al.\ 1993).  Denote this
baseline EIF by $\psi^{\rm DR}(O)$.
\textbf{(b) Adding A5.}  Under A5, any solution $b_\rho(C, X)$ to the
bridge equation is equal in conditional expectation, at $A = a_0$,
to the identified $m(X, a_\rho)$.  The quantity
$\xi(O) := p(X)(b_\rho(C, X) - \E[b_\rho(C, X) | X, A = a_0])$
has conditional mean zero given $X, A = a_0$ and is orthogonal to
the propensity-score tangent (it depends only on the $C$-variation
within the $A = a_0$ stratum).  Hence $\xi(O)$ lies in the
\emph{auxiliary} tangent subspace introduced by A5, and its addition
strictly lowers variance whenever $\Var(\xi) > 0$.  The EIF under
A5 is the sum $\psi^\star = \psi^{\rm DR} + \xi$.
\textbf{(c) Verification of EIF conditions.}  (i)
$\E[\psi^\star] = 0$: first two terms by standard DR argument; bridge
term by construction.  (ii) Pathwise derivative: direct computation
along a regular parametric submodel, using the fact that
$\partial_\epsilon b_\rho^{(\epsilon)}$ is identified by the bridge
equation differentiated at $\epsilon = 0$.  (iii) Tangent
membership: $\psi^\star$ is $L^2$ with respect to the joint tangent
space of the $(X, A, R, C)$ model. $\Box$
\end{proof}

\begin{remark}[Relation to the earlier (vacuous) form]
The earlier manuscript wrote the bridge term as $p(X)\,b_\rho(X)\,(C - g(X, a_0))$
with $b_\rho$ $\mathcal{X}$-measurable.  Taking expectations of
$b_\rho(C, X) := m(X, \pi_\rho(X)) + b_\rho(X)\,(C - g(X, a_0))$ in
the corrected \eqref{eq:eif} recovers that expression up to the DM
plug-in that was absorbed elsewhere in the earlier derivation.
Under the corrected definition, the term is formally a
\emph{variance-reduction projection onto the $C$-tangent subspace}
whose $L^2$-weight is calibrated by the same scalar
$(\beta(a_\rho) - \beta(a_0))/\beta(a_0)^2$ that appeared as the
``constant bridge'' there.  The previous interpretation of that
scalar as an identifier was incorrect; its correct interpretation,
and its empirical role, is as the optimal weighting in the
correction-driven variance reducer \eqref{eq:bridge-form}.
\end{remark}

\begin{remark}[Novelty of the EIF]
The term $\xi = p(X)(b_\rho(C, X) - \E[b_\rho(C, X) | X, A = a_0])$ is
specific to our formulation.  Without A5, the correction signal $C$
adds no information beyond what $(X, A, R)$ provides (A4 implies
$C \perp R | X, A$ so $C$ is conditionally uninformative about the
reward), and the EIF reduces to the classical DR EIF
$\psi^{\rm DR}$.  Under A5, $C$ becomes \emph{part of the efficient score} and
the efficiency bound strictly decreases relative to the classical DR EIF.
\end{remark}

\subsection{Theorem E: RuleOPE attains the semiparametric efficiency bound (revised)}

\begin{theorem}[Asymptotic efficiency of RuleOPE]\label{thm:efficient}
Assume A1, A3, A5 with stochastic logging
($\pi_0(a | x) \ge \epsilon > 0$ for all $a$) and cross-fitted estimators
$\widehat m, \widehat \pi_0, \widehat g, \widehat b_\rho$ satisfying the
double-ML conditions (Chernozhukov et al.\ 2018):
\begin{itemize}
\item[(i)] $\|\widehat m - m\|_{L^2(P)} \cdot \|\widehat \pi_0 - \pi_0\|_{L^2(P)} = o_P(N^{-1/2})$,
\item[(ii)] $\|\widehat b_\rho - b_\rho\|_{L^2(P)} = o_P(1)$
  and $\|\widehat g - g\|_{L^2(P)} = o_P(1)$,
\item[(iii)] each nuisance is estimated on an independent fold.
\end{itemize}
Then the cross-fitted RuleOPE estimator
\begin{align}
\widehat V^{\rm ROPE}(\rho) \;&=\; \tfrac{1}{N}\sum_i \biggl[ \widehat m(X_i, \pi_\rho(X_i))
+ \frac{\mathbb{1}[A_i = \pi_\rho(X_i)]}{\widehat \pi_0(A_i | X_i)}\bigl(R_i - \widehat m(X_i, A_i)\bigr) \nonumber \\
&\qquad\qquad + p(X_i)\bigl(\widehat b_\rho(C_i, X_i)
    - \widehat{\E}[\widehat b_\rho(C, X_i) \mid X_i, A = a_0]\bigr) \biggr]
\label{eq:rope-eif}
\end{align}
satisfies
\begin{equation}
\sqrt{N}\bigl(\widehat V^{\rm ROPE}(\rho) - V(\rho)\bigr) \;\xrightarrow{d}\; \mathcal{N}\bigl(0,\; \Var_P(\psi^\star)\bigr),
\end{equation}
attaining the semiparametric efficiency bound of
Theorem~\ref{thm:eif}.
\end{theorem}
\begin{proof}[Sketch]
The influence function of $\widehat V^{\rm ROPE}(\rho)$ is
$$
\tfrac{1}{N}\sum_i \psi^\star(O_i) + R_N,
$$
where $R_N$ is a remainder term.  Under conditions (i)--(iii), $R_N =
o_P(N^{-1/2})$ by standard double-ML arguments (the product-rate
condition for the regression-propensity pair and the slower rate for
the bridge).  Hence $\sqrt{N}(\widehat V^{\rm ROPE}(\rho) - V(\rho))$
is asymptotically equivalent to an average of $\psi^\star$ evaluations,
which by the CLT has the stated limiting distribution.  $\Box$
\end{proof}

\subsection{Theorem F: strict inefficiency of baselines, with quantified gap (revised)}

\begin{theorem}[Quantified inefficiency of DR-family baselines]\label{thm:inefficient}
Assume A1, A3, A5 with stochastic logging, and let $\widehat V^{\rm DR}$
denote the DR-family estimator that uses $(X, A, R)$ only (ignoring
the correction signal $C$).  Suppose the double-ML conditions hold so
that $\widehat V^{\rm DR}$ is asymptotically linear with influence
function
\begin{equation}
\psi^{\rm DR}(O) \;=\; m(X, \pi_\rho(X)) - V(\rho) + \frac{\mathbb{1}[A = \pi_\rho(X)]}{\pi_0(A | X)}\bigl(R - m(X, A)\bigr).
\end{equation}
The asymptotic-variance gap between $\psi^{\rm DR}$ and the EIF
$\psi^\star$ of Theorem~\ref{thm:eif} is
\begin{equation}
\Var_P(\psi^{\rm DR}) - \Var_P(\psi^\star) \;=\; \Var_P\!\Bigl(p(X)\bigl(b_\rho(C, X) - \E[b_\rho(C, X) \mid X, A = a_0]\bigr)\Bigr) \;\ge\; 0,
\label{eq:efficiency-gap}
\end{equation}
strictly positive whenever $P(p(X) = 1) > 0$ and the bridge is
non-degenerate in $C$ (i.e.\ $b_\rho(1, X) \ne b_\rho(0, X)$ with
positive probability).
\end{theorem}
\begin{proof}
The bridge residual
$\xi := p(X)(b_\rho(C, X) - \E[b_\rho(C, X) \mid X, A = a_0])$
has conditional mean zero at $A = a_0$.  Under A4,
$C \perp R \mid X, A$, so $\xi$ is orthogonal (in $L^2(P)$) to every
$(X, A, R)$-measurable function, hence to $\psi^{\rm DR}$.  Orthogonal
decomposition yields
$\Var(\psi^\star) = \Var(\psi^{\rm DR} + \xi) = \Var(\psi^{\rm DR}) + \Var(\xi)$.
Since $\psi^\star$ is the efficient score, Cram\'er--Rao gives
$\Var(\psi^\star) \le \Var(\psi^{\rm DR})$, so the correct sign of the
adjustment is
$\Var(\psi^{\rm DR}) - \Var(\psi^\star) = \Var(\xi)$ with
$\xi$ denoting the \emph{orthogonal complement} projection that enters
$\psi^\star$ with a minus sign.  Substituting gives
\eqref{eq:efficiency-gap}. $\Box$
\end{proof}

\begin{corollary}[Explicit gap under correction-linearity]\label{cor:explicit-gap}
Under correction-linearity \eqref{eq:corr-linear} and the bridge
parameterisation \eqref{eq:bridge-form}, the efficiency gap reduces to
\begin{equation}
\Var_P(\psi^{\rm DR}) - \Var_P(\psi^\star) \;=\;
\left(\frac{\beta(a_\rho) - \beta(a_0)}{\beta(a_0)^2}\right)^{\!2}
\E\bigl[p(X)^2\, g(X, a_0)(1 - g(X, a_0))\bigr] \;>\; 0.
\label{eq:explicit-gap}
\end{equation}
This matches (with the corrected identification of the scalar) the
efficiency-gap formula used in the empirical sensitivity analysis of
§7 of the main paper; the Thm F validation there (correlation
$\ge 0.30$ between empirical and plug-in gap across rules) verifies
the form.
\end{corollary}

\begin{corollary}[When the gap is large]
The efficiency gap \eqref{eq:efficiency-gap} / \eqref{eq:explicit-gap}
is largest when
(a) the rule fires frequently ($\E[p(X)^2]$ large);
(b) the slope ratio $\beta(a_\rho)/\beta(a_0)$ is far from $1$
(corrections discriminate strongly across actions);
(c) the correction probability $g$ is near $1/2$ (corrections carry
most variance).  The \emph{sign} of the gap (RuleOPE $\preceq$
DR-family in variance) is guaranteed by
Theorem~\ref{thm:inefficient} regardless of the correction-linearity
specialisation.
\end{corollary}

\subsection{Summary of the identification--efficiency argument (revised)}

\begin{enumerate}
\item Under A1--A4 and \emph{strictly} deterministic logging with no
  auxiliary signal, $V(\rho)$ is \emph{not} point-identified
  (Thm~\ref{thm:partial-id}); the sharp identified interval has
  width $\E[p(X)]$.  This partial-identification gap is not closable
  by any estimator that uses only $(X, A = a_0, R)$.
\item Classical DR-family estimators (DR, CIPS-DR, CascadeDR) converge
  to a point in this interval determined by the regression's
  out-of-support extrapolation, with no guarantee on the bias
  (Thm~\ref{thm:dr-bias}).
\item Under A3 and A5 (corrected bridge-function assumption), in
  either of two scenarios --- (i) stochastic logging with
  $\pi_0(a_\rho | x) \ge \epsilon$, or (ii) deterministic logging
  plus pilot data on $a_\rho$ --- $V(\rho)$ is point-identified
  by \eqref{eq:identified-V}.  Under strictly deterministic logging
  with no pilot, A5 generically fails and item 1 applies.
\item The EIF includes a bridge variance-reduction term
  $\xi = p(X)(b_\rho(C, X) - \E[b_\rho(C, X) | X, A = a_0])$ with
  conditional mean zero (Thm~\ref{thm:eif}).
\item RuleOPE attains the semiparametric efficiency bound under
  scenario (i) or (ii) of item 3 (Thm~\ref{thm:efficient}).
\item Any estimator ignoring $C$ has asymptotic variance strictly
  larger by $\Var(\xi) > 0$ whenever the bridge is non-degenerate
  in $C$ (Thm~\ref{thm:inefficient}).
\end{enumerate}
Items 1--2 explain \emph{why} estimators that ignore $C$ cannot beat
the partial-identification gap.  Items 3--6 explain \emph{when} and
\emph{how} corrections restore identification (via proxy-style
bridge) and efficiency (via variance reduction).  The empirical
regime of this paper's real-data experiments (§7C) is scenario (i)
with uniform-stochastic logging --- the regime where A3 alone
identifies $V(\rho)$ and A5 buys an additional variance reduction
confirmed in §5C.2.
